% \newpage
\subsection*{A4: Multi-Agent Coding}
\label{sec:app:optiguide}

In this subsection, we use \libName to build a multi-agent coding system based on OptiGuide~\cite{li2023large}, a system that excels at writing code to interpret optimization solutions and answer user questions, such as exploring the implications of changing a supply-chain decision or understanding why the optimizer made a particular choice. 
The second sub-figure of Figure~\ref{fig:app_summary} shows the \libName-based implementation.  The workflow is as follows:
the end user sends questions, such as ``\emph{What if we prohibit shipping from supplier 1 to roastery 2?}" to the Commander agent.  
The Commander coordinates with two assistant agents, including the Writer and the Safeguard, to answer the question. The Writer will craft code 
and send the code to the Commander. 
After receiving the code, the Commander checks the code safety with the Safeguard; if cleared, the Commander will use external tools (e.g., Python) to execute the code, and request the Writer to interpret the execution results. For instance, the writer may say ``\emph{if we prohibit shipping from supplier 1 to roastery 2, the total cost would increase by 10.5\%.}" The Commander then provides this concluding answer to the end user.  If, at a particular step, there is an exception, e.g., security red flag raised by Safeguard, the Commander redirects the issue back to the Writer with debugging information.
The process might be repeated multiple times until the user's question is answered or timed-out. 

With \libName the core workflow code for OptiGuide was reduced from over 430 lines to 100 lines, leading to significant productivity improvement. We provide a detailed comparison of user experience with ChatGPT+Code Interpreter and \libName-based OptiGuide in Appendix~\ref{appendix:app:optiguide}, where we show that \libName-based OptiGuide could save around 3x of user's time and reduce user interactions by 3 - 5 times on average.  We also conduct an ablation showing that multi-agent abstraction is necessary. Specifically, we construct a single-agent approach where a single agent conducts both the code-writing and safeguard processes. We tested the single- and multi-agent approaches on a dataset of 100 coding tasks, which is crafted to include equal numbers of safe and unsafe tasks.
Evaluation results as reported in Figure \ref{fig:res:sageguard}
show that the multi-agent design boosts the F-1 score in identifying unsafe code by 8\% (with GPT-4) and 35\% (with GPT-3.5-turbo). 

